\documentclass[10pt,a4paper]{amsart}
\usepackage[margin=19mm]{geometry}
\usepackage{amsmath,amssymb,mathtools}
\usepackage[scr=rsfs,cal=euler]{mathalfa}
\usepackage{xfrac}
\usepackage[unicode,hidelinks,bookmarks=false]{hyperref}
\newcommand{\ie}{i.e.\ }
\newcommand{\cf}{cf.\ }
\newcommand{\resp}{respectively\ }
\newcommand{\Subsection}{\subsection}
\providecommand{\nobreakdash}{\nobreak\hbox{-}}
\setlength{\emergencystretch}{2em}
\multlinegap0pt
\usepackage{mdframed}
\usepackage{mdframed}
\usepackage{mdframed}
\usepackage{mdframed}
\usepackage{mdframed}
\usepackage{bm}
\usepackage{dsfont}
\usepackage{stackengine}
\usepackage{cancel}
\usepackage{xcolor}
\usepackage{tcolorbox}
\usepackage{algorithm}\usepackage{algpseudocode}
\usepackage{tikz}
\usepackage[shortlabels]{enumitem}
\newtheorem{lemma}{Lemma}
\newcommand\abs[1]{\left\lvert#1\right\rvert}
\newcommand\absbig[1]{\big\lvert#1\big\rvert}
\newcommand{\bG}{{\bm{G}}}
\newcommand{\bS}{{\bm{S}}}
\newcommand{\bX}{{\bm{X}}}
\newcommand{\bbeta}{{\boldsymbol\beta}}
\newcommand{\bone}{{\bm{1}}}
\newcommand{\bs}{{\bm{s}}}
\newcommand{\bu}{{\bm{u}}}
\newcommand{\bv}{{\bm{v}}}
\newcommand{\bw}{{\bm{w}}}
\newcommand{\bx}{{\bm{x}}}
\newcommand{\by}{{\bm{y}}}
\newcommand{\bzero}{\boldsymbol{0}}
\newcommand\comple[1]{#1^c}
\newcommand{\diag}{\mathrm{diag}}
\newcommand\mathopmin[1]{\mathop{\min}_{#1}}
\newcommand{\mdframecolor}{gray!10}
\newcommand{\mdframehideline}{true}
\newcommand\normone[1]{\left\lVert#1\right\rVert_1}
\newcommand\normonebig[1]{\big\lVert#1\big\rVert_1}
\newcommand\normtwo[1]{\left\lVert#1\right\rVert_2}
\newcommand{\real}{\mathbb{R}}
\def\sM{{\mathbb{M}}}
\def\sS{{\mathbb{S}}}
\newcommand{\sgn}{\mathrm{sign}}
\DeclareMathOperator{\sign}{sign}
\newcommand{\widehatbbeta}{{\widehat\bbeta}}
\newcommand{\widehatbeta}{{\widehat\beta}}
\newcommand{\widehatby}{\widehat{\bm{y}}}
\newcommand{\widehatc}{\widehat{c}}
\title{A First Course in Sparse Optimization}
\author{Jun Lu}
\date{Source: arXiv:2601.06173v1 (CC BY 4.0)}
\begin{document}
\maketitle

\noindent\emph{Source and licence.} A First Course in Sparse Optimization. Version arXiv:2601.06173v1 (CC BY 4.0), distributed by arXiv under the Creative Commons Attribution 4.0 International licence, http://creativecommons.org/licenses/by/4.0/, as recorded in the arXiv licence metadata for this exact version and shown on the abstract page https://arxiv.org/abs/2601.06173v1. Full licence notice: \texttt{licenses/arxiv-2601.06173-CC-BY-4.0.txt}.\par\smallskip

\noindent\textit{Editorial scope.} Counted unit: the article's lemma lemma:larslasso\_ud\_est (source0.tex lines 23348--23361). The article's complete original proof of it is delivered with it (source0.tex lines 23364--23467, one \texttt{proof} environment of 104 lines, which the article places away from the statement). No mathematics is written, repaired or re-derived: the statement and the proof are the article's own lines, in the article's own notation, and no retained line is substituted or re-flowed.  Retained uncounted as the source-local context the proof needs: lines 10881--10886; lines 23114--23116; lines 23165--23167.  Nothing was omitted from any retained span, so the package carries no omission record.  No retained line contains a citation, so the wrapper carries no bibliography and no bibliography entry.  No retained line contains a figure environment, an image-inclusion command or drawing code, so no image is delivered and the package declares no asset dependency.  Editorial formatting only, in addition: an AMS \texttt{amsart} preamble that restates the article's own declarations and macros used by the retained lines, namely the environments \texttt{lemma} on separate counters, together with the standard packages those lines need.  The retained environments print in this document as Lemma 1 in order of appearance, whereas the article prints the same items with the numbering of its own full document.  The complete native source of the article as distributed in the arXiv e-print for this exact version is bundled unchanged as \texttt{sources/192393/source0.tex}.
\begin{align}
(\textbf{Constrained LASSO}):\qquad &\mathopmin{\normone{\bbeta}\leq \Sigma}f(\bbeta)\triangleq \frac{1}{2}\normtwo{\by - \bX \bbeta}^2;
\label{opt:lc} \tag{L$_{C}$}\\
(\textbf{Lagrangian LASSO}):\qquad &\mathopmin{\bbeta \in \real^p}F(\bbeta)\triangleq \left\{ \frac{1}{2} \normtwo{\by - \bX \bbeta}^2 + \lambda \normone{\bbeta} \right\}. 
\label{opt:ll} \tag{L$_{L}$}
\end{align}

\begin{equation}\label{equation:ang_bisec_eq0}
\bX_t \triangleq  [\ldots, s_i \bx_i, \ldots]_{i \in \sS},
\end{equation}

\begin{equation}\label{equation:lars_upd2}
s_i \triangleq \sgn\{\widehatc_i\}, \quad \text{for } i \in \sS,
\end{equation}

\begin{lemma}[Piecewise linear update of the fitted values in LARS--LASSO]\label{lemma:larslasso_ud_est}
For $\Sigma \in \sM$, the LASSO estimates $\widehatby(\Sigma)$ satisfy
\begin{equation}\label{equation:haty_sigma_linear}
\widehatby(\Sigma) = \widehatby(\Sigma_0) + \sigma_\sS (\Sigma - \Sigma_0) \bu_\sS
\end{equation}
where $\bu_\sS$ is the equiangular vector $\bu_{\sS}=\bX_\sS \bw_\sS$, $\bw_\sS = \sigma_\sS \bG_\sS^{-1} \bone_\sS$.
Equivalently, in terms of the nonzero regression coefficients $\widehatbbeta_\sS(\Sigma)$,
\begin{equation}\label{equation:beta_sigma_linear}
\widehatbbeta_\sS(\Sigma) = \widehatbbeta_\sS(\Sigma_0) + \bS_\sS \sigma_\sS (\Sigma - \Sigma_0) \bw_\sS,
\end{equation}
where $\bS_\sS \triangleq  \diag(\bs_\sS)$ is the diagonal matrix with diagonal elements $s_i = \sign(\widehatbeta_i)$, $i \in \sS$, and $\bs_\sS$ is a vector containing these sign elements.~\footnote{
The matrix $\bS_\sS$ appears because, by definitions~\eqref{equation:ang_bisec_eq0} and~\eqref{equation:lars_upd2}, we have  $\widehatby(\Sigma) = \bX \widehatbbeta(\Sigma) = \bX_\sS \bS_\sS \widehatbbeta_\sS(\Sigma)$.
}
\end{lemma}

\begin{proof}[of Lemma~\ref{lemma:larslasso_ud_est}]
The lemma says that, for $\Sigma$ in $\sM$ where the active set $\sS$ and signs $\bs_{\sS} $
remain fixed, $\widehatby(\Sigma)$ moves linearly along the equiangular vector $\bu_\sS$ determined by $\sS$. 
Since $\widehatbbeta(\Sigma)$ satisfies the constrained LASSO problem \eqref{opt:lc} and has support $\sS$, it also minimizes
\begin{equation*}
f(\widehatbbeta_\sS) = \frac{1}{2}\normtwo{\by - \bX_\sS \bS_\sS \widehatbbeta_\sS}^2
\end{equation*}
subject to
\begin{equation}\label{eq:constraint}
\sum_{i \in \sS} s_i \widehatbeta_i = \Sigma 
\qquad \text{and} \qquad 
\sign(\widehatbeta_i) = s_i \quad \text{for } i \in \sS.
\end{equation}
The inequality constraint in the constraint LASSO problem \eqref{opt:lc} can be replaced by $\normonebig{\widehatbbeta} = \Sigma$ as long as $\Sigma$ is less than $\sum \absbig{\widehatbeta_i}$ for the full $p$-variable OLS solution $\widehatbbeta_p$.
Moreover, the fact that the minimizing point $\widehatbbeta_\sS(\Sigma)$ occurs strictly inside the simplex~\eqref{eq:constraint}, combined with the strict convexity of $f(\widehatbbeta_\sS)$, implies we can drop the second condition in~\eqref{eq:constraint} so that $\widehatbbeta_\sS(\Sigma)$ solves
\begin{equation}\label{eq:min_pieceproblem}
\min \quad \{ f(\widehatbbeta_\sS) \} \quad \text{s.t.} \quad \sum_{i \in \sS} s_i \widehatbeta_i = \Sigma.
\end{equation}
Introducing a Lagrange multiplier $\lambda$, \eqref{eq:min_pieceproblem} becomes
$$
\min \quad \frac{1}{2} \normtwo{\by - \bX_\sS \bS_\sS \widehatbbeta_\sS}^2 + \lambda \sum_{i \in \sS} s_i \widehatbeta_i.
$$
Differentiating with respect to $\bbeta_\sS$ and setting the gradient to zero yields
\begin{align}
&- \bS_\sS \bX_\sS^\top (\by - \bX_\sS \bS_\sS \widehatbbeta_\sS) + \lambda \bS_\sS \bone_\sS = \bzero \label{eq:larslasso_stat_condition010} \\
&\quad \implies\quad
\bbeta_\sS = (\bS_\sS\bX_\sS^\top \bX_\sS\bS_\sS)^{-1}(\bS_\sS\bX_\sS^\top \by - \lambda \bs_\sS). \label{eq:larslasso_stat_condition}
\end{align}






Now consider two values $\Sigma_1, \Sigma_2\in \sM$ with $\Sigma_0 < \Sigma_1 < \Sigma_2$, 
and let $\lambda_1 > \lambda_2$ be the corresponding Lagrange multipliers  (since larger $\Sigma$ implies smaller penalty weight).  
Subtracting the optimality conditions for $\Sigma_1$ and $\Sigma_2$ gives
$$
\bX_\sS^\top \bX_\sS \bS_\sS \big(\widehatbbeta_\sS(\Sigma_2) - \widehatbbeta_\sS(\Sigma_1)\big) = (\lambda_1 - \lambda_2) \bone_\sS,
$$
which implies
$$
\widehatbbeta_\sS(\Sigma_2) - \widehatbbeta_\sS(\Sigma_1) = (\lambda_1 - \lambda_2) \bS_\sS \bG_\sS^{-1} \bone_\sS.
$$
However, we can also prove $\bs_\sS^\top [\widehatbbeta_\sS(\Sigma_2) - \widehatbbeta_\sS(\Sigma_1)] = \Sigma_2 - \Sigma_1$ according to the LASSO definition:

\begin{mdframed}[hidealllines=\mdframehideline,backgroundcolor=\mdframecolor,frametitle={Relationship between $\widehatbbeta_\sS(\Sigma_2) - \widehatbbeta_\sS(\Sigma_1) $ and $ \Sigma_2 - \Sigma_1$:}]
For the active set, $\Sigma = \normone{\bbeta} = \sum_{i \in \sS} \abs{\beta_i} = \bs_\sS^\top \bbeta_\sS$.
Substitute \eqref{eq:larslasso_stat_condition}:
$$
\Sigma = \bs_\sS^\top (\bS_\sS\bX_\sS^\top \bX_\sS\bS_\sS)^{-1}(\bS_\sS\bX_\sS^\top \by - \lambda \bs_\sS).
$$
Define
$a \triangleq \bs_\sS^\top (\bS_\sS\bX_\sS^\top \bX_\sS\bS_\sS)^{-1} \bS_\sS\bX_\sS^\top \by$ 
and  $b \triangleq \bs_\sS^\top (\bS_\sS\bX_\sS^\top \bX_\sS\bS_\sS)^{-1} \bs_\sS$.
Then
$$
\Sigma = a - \lambda b
\quad \Longrightarrow \quad
\lambda = \frac{a - \Sigma}{b}.
$$
Thus, $\lambda$ is an affine (linear) function of $\Sigma$
while $\sS$ and $\bs_\sS$ remain fixed.
Differentiate \eqref{eq:larslasso_stat_condition} with respect to $\Sigma$:
$
\frac{\partial \bbeta_\sS}{\partial \Sigma}
= - (\bS_\sS\bX_\sS^\top \bX_\sS\bS_\sS)^{-1} \bs_\sS \cdot \frac{\partial \lambda}{\partial \Sigma}
$.
From $\lambda = (a - \Sigma)/b$, we get $\dfrac{\partial \lambda}{\partial \Sigma} = -\dfrac{1}{b}$.
Hence
$$
\frac{\partial\bbeta_\sS}{\partial \Sigma}
= \frac{(\bS_\sS\bX_\sS^\top \bX_\sS\bS_\sS)^{-1} \bs_\sS}
{\bs_\sS^\top (\bS_\sS\bX_\sS^\top \bX_\sS\bS_\sS)^{-1} \bs_\sS}
, \qquad
\frac{\partial\bbeta_{\comple{\sS}}}{\partial\Sigma} = \bzero.
$$
If $\Sigma_1$ and $\Sigma_2$ lie within the same region
where the active set $\sS$ and sign vector $\bs_\sS$ remain unchanged, then
\begin{equation}
\bbeta_2 - \bbeta_1
= (\Sigma_2 - \Sigma_1)
\frac{(\bS_\sS\bX_\sS^\top \bX_\sS\bS_\sS)^{-1} \bs_\sS}
{\bs_\sS^\top (\bS_\sS\bX_\sS^\top \bX_\sS\bS_\sS)^{-1} \bs_\sS}.
\end{equation}
That is, $\bbeta$ changes linearly with $\Sigma$
along the fixed direction
$
\bv \triangleq \frac{(\bX_\sS^\top \bX_\sS)^{-1} \bs_\sS}
{\bs_\sS^\top (\bX_\sS^\top \bX_\sS)^{-1} \bs_\sS}
$.
\end{mdframed}
Therefore,
$$
\Sigma_2 - \Sigma_1 = (\lambda_1 - \lambda_2) \bs_\sS^\top \bS_\sS \bG_\sS^{-1} \bone_\sS = (\lambda_1 - \lambda_2) \bone_\sS^\top \bG_\sS^{-1} \bone_\sS = (\lambda_1 - \lambda_2) \sigma_\sS^{-2},
$$
and hence
\begin{equation}\label{eq:beta_diff_final}
\widehatbbeta_\sS(\Sigma_2) - \widehatbbeta_\sS(\Sigma_1) 
= \bS_\sS \sigma_\sS^2 (\Sigma_2 - \Sigma_1) \bG_\sS^{-1} \bone_\sS 
= \bS_\sS \sigma_\sS (\Sigma_2 - \Sigma_1) \bw_\sS.
\end{equation}
Invoking $\Sigma_2 = \Sigma$ and $\Sigma_1 \to \Sigma_0$ gives~\eqref{equation:beta_sigma_linear} by the continuity of $\widehatbbeta(\Sigma)$, and implying the desired result \eqref{equation:haty_sigma_linear}. 
\end{proof}

\end{document}
